\documentclass[12pt,a4paper]{article}

\usepackage[margin=1in]{geometry}
\usepackage{graphicx}
\usepackage{amsmath}
\usepackage{booktabs}
\usepackage{array}
\usepackage{float}
\usepackage{xcolor}
\usepackage{hyperref}

\title{
    A Multi-Agent Search-Based Network Attack and Defense Simulation
}

\author{
    Fundamentals of Artificial Intelligence\\
    23CSE401
}

\date{\today}

\begin{document}

\maketitle

\begin{abstract}

This project presents a multi-agent artificial intelligence simulation of a
network attack and defense scenario. The system models an attacker attempting
to reach a critical database through a simulated computer network while a
defender attempts to prevent the attack by patching, isolating, and monitoring
hosts.

The project applies classical AI search techniques including A* search,
Minimax, and Alpha-Beta pruning. Random and greedy strategies are also used as
baselines. The system evaluates attack paths, models competing attacker and
defender decisions, demonstrates the resulting search tree, and compares the
computational cost of Minimax and Alpha-Beta pruning.

The experiments show that Alpha-Beta pruning can significantly reduce the
number of explored search nodes while producing the same decision and
evaluation score as Minimax.

\end{abstract}


\tableofcontents

\newpage


% ============================================================
\section{Introduction}
% ============================================================

Artificial intelligence search techniques can be used to model decision
making in environments where an agent must select actions from a set of
possible alternatives.

This case study applies multi-agent search to a simulated network security
scenario. An attacker attempts to move through a network and reach a critical
database, while a defender attempts to prevent this from happening.

The problem is modeled as a competitive, turn-based game. The attacker acts
as the maximizing player and the defender acts as the minimizing player.
A* search is used for attack-path planning, while Minimax and Alpha-Beta
pruning are used to model strategic interaction between the two agents.

The implementation is entirely simulation-based. It does not perform attacks
against real systems.


% ============================================================
\section{Problem Domain and Problem Statement}
% ============================================================

\subsection{Problem Domain}

The problem domain is network attack and defense.

The simulated network contains multiple interconnected hosts with different
values and vulnerabilities. The attacker starts from the Internet and attempts
to reach the database. The defender attempts to prevent the attacker from
reaching the target.

\subsection{Problem Statement}

The objective is to design a multi-agent AI system in which an attacker and a
defender interact within a simulated network.

The attacker must select actions that help it move toward the target database,
while the defender selects actions that make the attack more difficult.

The system should:

\begin{itemize}
    \item represent the network as a weighted graph;
    \item model attacker and defender actions;
    \item find attack paths using A* search;
    \item model competing decisions using Minimax;
    \item improve Minimax efficiency using Alpha-Beta pruning;
    \item compare intelligent strategies with simpler baselines;
    \item demonstrate normal and difficult cases; and
    \item analyze search complexity.
\end{itemize}


% ============================================================
\section{System Description}
% ============================================================

The network is represented as a weighted graph.

The simulated hosts are:

\begin{itemize}
    \item Internet
    \item Workstation
    \item Web Server
    \item Application Server
    \item Database
    \item Admin PC
    \item File Server
    \item Backup Server
\end{itemize}

Each host contains information such as its type, value, vulnerabilities,
patch status, compromise status, monitoring status, and isolation status.

Connections between hosts have different traversal costs.

The database is treated as the critical target of the attacker.


% ============================================================
\section{Agents}
% ============================================================

\subsection{Attacker Agent}

The attacker starts at the Internet and attempts to reach the database.

The attacker can:

\begin{itemize}
    \item exploit a vulnerable connected host;
    \item move through the network as hosts are compromised; and
    \item wait.
\end{itemize}

The attacker's objective is to maximize its evaluation score and ultimately
reach the target.

\subsection{Defender Agent}

The defender attempts to prevent the attacker from reaching the database.

The defender can:

\begin{itemize}
    \item patch vulnerable hosts;
    \item isolate hosts;
    \item monitor hosts; and
    \item wait.
\end{itemize}

The defender attempts to minimize the attacker's evaluation score.


% ============================================================
\section{PEAS Description}
% ============================================================

\subsection{Performance Measure}

For the attacker, the performance measure includes:

\begin{itemize}
    \item reaching the target database;
    \item reducing the cost of the attack path; and
    \item compromising valuable hosts.
\end{itemize}

For the defender, the performance measure includes:

\begin{itemize}
    \item preventing access to the database;
    \item increasing the attacker's remaining path cost; and
    \item reducing available attack routes.
\end{itemize}

\subsection{Environment}

The environment is a simulated computer network represented as a weighted
graph. It contains hosts, network connections, vulnerabilities, patch states,
isolation states, monitoring states, and compromised hosts.

\subsection{Actuators}

The attacker actuators are:

\begin{itemize}
    \item Exploit
    \item Wait
\end{itemize}

The defender actuators are:

\begin{itemize}
    \item Patch
    \item Isolate
    \item Monitor
    \item Wait
\end{itemize}

\subsection{Sensors}

The agents use information from the simulated environment, including:

\begin{itemize}
    \item current attacker position;
    \item network connections;
    \item host vulnerabilities;
    \item patch status;
    \item isolation status;
    \item compromised hosts; and
    \item target location.
\end{itemize}


% ============================================================
\section{Environment Properties}
% ============================================================

The environment is characterized as follows.

\begin{table}[H]
\centering
\begin{tabular}{p{0.25\textwidth} p{0.65\textwidth}}
\toprule
\textbf{Property} & \textbf{Classification and Justification} \\
\midrule

Observability &
Fully observable. The search algorithm has access to the simulated network
state, including host vulnerabilities, patches, isolation, and compromise
information. \\

Determinism &
Deterministic. A valid action produces a predictable next state. \\

Episodic / Sequential &
Sequential. Current actions modify the network state and affect future
decisions. \\

Static / Dynamic &
Dynamic. The state of the network changes as the attacker and defender take
actions. \\

Discrete / Continuous &
Discrete. Hosts, actions, turns, and game states are represented using
discrete values. \\

Single-agent / Multi-agent &
Multi-agent. An attacker and a defender interact in the same environment. \\

Cooperative / Competitive &
Competitive. The attacker attempts to reach the target while the defender
attempts to prevent this. \\

\bottomrule
\end{tabular}
\caption{Environment properties}
\end{table}


% ============================================================
\section{Agent Type}
% ============================================================

The attacker and defender are modeled as goal-based, utility-driven agents.

The attacker selects actions based on their expected effect on its evaluation
score and its ability to reach the database.

The defender selects actions based on their effect on reducing the attacker's
evaluation score and restricting possible attack paths.

The agents are competitive rather than cooperative because their objectives
conflict.

The Minimax formulation treats the attacker as the MAX player and the
defender as the MIN player.

This is a simplified turn-based model of adversarial decision making rather
than a claim that real-world cyber attackers and defenders always behave as
strictly alternating zero-sum players.


% ============================================================
\section{State Representation}
% ============================================================

A game state contains:

\begin{itemize}
    \item attacker position;
    \item target host;
    \item compromised hosts;
    \item patched hosts;
    \item isolated hosts;
    \item current turn;
    \item current search depth; and
    \item the network graph.
\end{itemize}

The state is modified whenever an attacker or defender action is applied.

The target is considered reached when the database appears in the set of
compromised hosts.


% ============================================================
\section{Search Techniques}
% ============================================================

\subsection{A* Search}

A* search is used to determine a low-cost feasible attack path from the
attacker's current position to the target.

A* uses:

\begin{equation}
f(n) = g(n) + h(n)
\end{equation}

where $g(n)$ is the cost accumulated from the starting state and $h(n)$ is an
estimate of the remaining cost.

The implementation uses the weighted shortest-path distance to the target as
the heuristic. The path generation also respects security restrictions such
as patched and isolated hosts.

One observed initial search produced:

\begin{verbatim}
Internet -> Web Server -> Application Server -> Database
\end{verbatim}

with total path cost:

\[
1 + 1 + 2 = 4.
\]

\subsection{Minimax}

Minimax models the interaction between the attacker and defender.

The attacker is the MAX player and attempts to maximize the evaluation score.

The defender is the MIN player and attempts to minimize the evaluation score.

The algorithm recursively explores possible future actions up to a specified
search depth.

\subsection{Alpha-Beta Pruning}

Alpha-Beta pruning optimizes Minimax by removing branches that cannot affect
the final decision.

The algorithm maintains two bounds:

\begin{itemize}
    \item $\alpha$: the best value currently guaranteed to MAX;
    \item $\beta$: the best value currently guaranteed to MIN.
\end{itemize}

A branch can be pruned when:

\begin{equation}
\beta \leq \alpha
\end{equation}

Alpha-Beta should produce the same decision and evaluation score as Minimax
when both algorithms use the same depth and evaluation function.


% ============================================================
\section{Evaluation Function}
% ============================================================

The evaluation function evaluates a state from the attacker's perspective.

A higher score represents a better state for the attacker.

The evaluation considers:

\begin{itemize}
    \item whether the target has been reached;
    \item the A* attack-path cost;
    \item the values of compromised hosts;
    \item the number of isolated hosts; and
    \item the number of patched hosts.
\end{itemize}

The simplified evaluation is based on the following components:

\begin{equation}
Score =
PathScore +
CompromisedValue -
IsolationPenalty -
PatchPenalty
\end{equation}

Reaching the target receives a very large positive score, while reaching the
maximum search depth without success receives a very low score.


% ============================================================
\section{Search Tree and State Space}
% ============================================================

The search tree represents possible sequences of attacker and defender
actions.

The initial state is:

\begin{itemize}
    \item Attacker position: Internet
    \item Target: Database
    \item Turn: Attacker
\end{itemize}

At the first level, the attacker can choose actions such as:

\begin{itemize}
    \item Exploit Web Server
    \item Exploit Workstation
    \item Wait
\end{itemize}

The defender then generates possible responses such as:

\begin{itemize}
    \item Patch a vulnerable host
    \item Isolate a host
    \item Monitor a host
    \item Wait
\end{itemize}

The implemented search-tree demonstration was executed to depth 2 and
generated these attacker branches and their corresponding defender branches.

The search tree demonstrates that the number of possible states grows
rapidly because each agent has multiple available actions.


% ============================================================
\section{Working Cases and Edge Cases}
% ============================================================

\subsection{Working Case 1: Normal Attack}

The attacker follows a valid route:

\begin{verbatim}
Internet -> Web Server -> Application Server -> Database
\end{verbatim}

The target is successfully reached.

Experimental result:

\begin{verbatim}
Working Case 1: True
\end{verbatim}

\subsection{Working Case 2: Defender Intervention}

The defender patches the Admin PC before the attacker can use that route.

The attacker can still attempt another route, demonstrating that defensive
actions modify the available attack space.

Experimental result:

\begin{verbatim}
Working Case 2: False
\end{verbatim}

The database is not reached in this scenario.

\subsection{Edge Case 1: Unreachable Target}

The defender isolates both entry routes from the Internet:

\begin{itemize}
    \item Web Server
    \item Workstation
\end{itemize}

The attacker therefore has no feasible route into the internal network.

Experimental result:

\begin{verbatim}
Edge Case 1: False
\end{verbatim}

\subsection{Edge Case 2: Maximum Search Depth}

The maximum depth is deliberately reduced to two actions.

The attacker and defender both wait, causing the search to reach its maximum
depth before the target is reached.

Experimental result:

\begin{verbatim}
Edge Case 2: depth=2 terminal=True
\end{verbatim}


% ============================================================
\section{Complexity Analysis}
% ============================================================

Let $b$ represent the average branching factor and $d$ represent the search
depth.

For standard Minimax, the approximate time complexity is:

\begin{equation}
O(b^d)
\end{equation}

The depth-first implementation requires approximately:

\begin{equation}
O(bd)
\end{equation}

space.

For Alpha-Beta pruning, the worst-case time complexity remains:

\begin{equation}
O(b^d)
\end{equation}

However, with favorable move ordering, Alpha-Beta can approach:

\begin{equation}
O(b^{d/2})
\end{equation}

The experiment demonstrates the practical reduction in explored states as
the search depth increases.


% ============================================================
\section{Experimental Results}
% ============================================================

Minimax and Alpha-Beta were evaluated at search depths from 1 to 4.

\begin{table}[H]
\centering
\begin{tabular}{ccccc}
\toprule
\textbf{Depth} &
\textbf{Minimax Nodes} &
\textbf{Alpha-Beta Nodes} &
\textbf{Nodes Saved} &
\textbf{Reduction} \\
\midrule
1 & 4 & 4 & 0 & 0.00\% \\
2 & 49 & 35 & 14 & 28.57\% \\
3 & 159 & 78 & 81 & 50.94\% \\
4 & 1707 & 591 & 1116 & 65.38\% \\
\bottomrule
\end{tabular}
\caption{Minimax versus Alpha-Beta search complexity}
\end{table}

At depth 4:

\begin{itemize}
    \item Minimax explored 1707 nodes.
    \item Alpha-Beta explored 591 nodes.
    \item 1116 nodes were saved.
    \item The reduction was 65.38\%.
    \item Both algorithms selected the same action.
    \item Both algorithms produced the same evaluation score.
\end{itemize}

This demonstrates that Alpha-Beta pruning can substantially reduce search
effort without changing the resulting decision.


% ============================================================
\section{Strategy Comparison}
% ============================================================

Four attacker strategies were compared:

\begin{enumerate}
    \item Random attacker
    \item Greedy/A* attacker
    \item Minimax attacker
    \item Alpha-Beta attacker
\end{enumerate}

One experimental run produced the following results.

\begin{table}[H]
\centering
\begin{tabular}{p{0.27\textwidth} p{0.28\textwidth} p{0.18\textwidth} p{0.18\textwidth}}
\toprule
\textbf{Strategy} & \textbf{Action} & \textbf{Score} & \textbf{Nodes} \\
\midrule
Random &
Exploit Web Server &
75 &
-- \\

Greedy/A* &
Exploit Workstation &
82 &
-- \\

Minimax &
Exploit Workstation &
79 &
1707 \\

Alpha-Beta &
Exploit Workstation &
79 &
591 \\

\bottomrule
\end{tabular}
\caption{Attacker strategy comparison}
\end{table}

The random strategy is nondeterministic, so its selected action and score can
change between runs.

Minimax and Alpha-Beta selected the same action and produced the same score.
Alpha-Beta achieved this result while exploring substantially fewer nodes.


% ============================================================
\section{Comparison with Alternative Search Technique}
% ============================================================

The primary comparison is between Minimax and Alpha-Beta pruning.

Minimax explores the game tree without pruning. As the depth increases, the
number of states grows rapidly.

Alpha-Beta pruning produces the same optimal decision under the same
evaluation and depth but avoids branches that cannot influence the final
decision.

The experimental results show:

\begin{itemize}
    \item 0\% node reduction at depth 1;
    \item 28.57\% reduction at depth 2;
    \item 50.94\% reduction at depth 3; and
    \item 65.38\% reduction at depth 4.
\end{itemize}

Therefore, Alpha-Beta becomes increasingly beneficial as the search space
grows.

A* is also used as a complementary search technique. Unlike Minimax, A* is
used for single-agent path planning and does not model the defender's future
decisions. The Greedy/A* attacker therefore provides a useful baseline for
comparison with the adversarial search strategies.


% ============================================================
\section{Limitations}
% ============================================================

The simulation is an abstract model rather than a real-world cyber range.

The main limitations are:

\begin{itemize}
    \item The environment is fully observable.
    \item Actions are deterministic.
    \item The attacker and defender take alternating turns.
    \item Network behavior is simplified.
    \item Real vulnerabilities and exploitation mechanisms are abstracted.
    \item Detection uncertainty and stochastic behavior are not modeled.
    \item The defender cannot directly patch or isolate the critical database
    target.
\end{itemize}

The final limitation is intentional. The database is treated as the critical
asset, so the defender must protect it indirectly by securing intermediate
hosts and routes.

The Minimax formulation should therefore be interpreted as a modeled,
turn-based competitive baseline rather than a claim that real cyber
attacks occur as strict two-player zero-sum games.


% ============================================================
\section{Conclusion}
% ============================================================

This project demonstrates the application of classical artificial
intelligence search techniques to a multi-agent network attack and defense
simulation.

A* search provides attack-path planning, while Minimax models the strategic
interaction between an attacker and defender. Alpha-Beta pruning improves the
efficiency of Minimax by eliminating branches that cannot affect the final
decision.

The experiments demonstrate that the search space grows rapidly with depth.
At depth 4, Alpha-Beta reduced the number of explored nodes by 65.38\% while
producing the same action and score as Minimax.

The project therefore demonstrates multi-agent decision making, graph-based
search, adversarial search, search-tree generation, edge-case handling, and
complexity analysis within a reproducible simulation environment.


\end{document}